% Part: first-order-logic
% Chapter: completeness
% Section: complete-sets

% Definition of complete consistent sets. Properties of complete sets required
% for completeness proved are in provability.tex in the
% chapter on the proof system used.

\documentclass[../../../include/open-logic-section]{subfiles}

\begin{document}

\iftag{FOL}
      {\olfileid{fol}{com}{ccs}}
      {\olfileid{pl}{com}{ccs}}
      
\olsection{Volledige consistente verzamelingen van \usetoken{P}{sentence}}

\begin{defn}[Volledige verzameling]
\ollabel{def:complete-set} Een verzameling~$\Gamma$ van !!{sentence}s is
\emph{!!{complete}} dan en slechts dan als voor elke !!{sentence}~$!A$
geldt dat $!A \in \Gamma$ of $\lnot !A \in \Gamma$.
\end{defn}

\begin{explain}
!!^{complete} verzamelingen zinnen laten geen vragen onbeantwoord. Voor
elke !!{sentence}~$!A$ ``zegt'' $\Gamma$ of $!A$ waar dan wel onwaar is.
Het belang van !!{complete} verzamelingen reikt verder dan het bewijs van
de volledigheidsstelling. Een theorie die !!{complete} en axiomatiseerbaar
is, is bijvoorbeeld altijd beslisbaar.
\end{explain}

\begin{explain}
!!^{complete} consistente verzamelingen zijn van belang voor het
volledigheidsbewijs, omdat we kunnen garanderen dat elke consistente
verzameling !!{sentence}s~$\Gamma$ opgenomen is in !!a{complete} consistente
verzameling~$\Gamma^*$. !!^a{complete} consistente verzameling bevat voor
elke !!{sentence}~$!A$ ofwel $!A$ ofwel haar ontkenning $\lnot !A$, maar
niet beide. Dit geldt in het bijzonder voor \iftag{FOL}{atomaire
!!{sentence}s}{propositionele variabelen}. Uit !!a{complete} consistente
verzameling\iftag{FOL}{ in een taal die op passende wijze met
!!{constant}s is uitgebreid}{} kunnen we daarom
\iftag{FOL}{!!a{structure}}{!!a{valuation}} construeren waarin de
\iftag{FOL}{interpretatie van !!{predicate}s}{waarheidswaarde die aan
propositionele variabelen wordt toegekend} wordt bepaald door de
\iftag{FOL}{atomaire !!{sentence}s}{!!{propositional variable}s} die
tot~$\Gamma^*$ behoren. Vervolgens kan worden aangetoond dat deze
\iftag{FOL}{!!{structure}}{!!{valuation}} alle !!{sentence}s in~$\Gamma^*$
waar maakt, en dus ook alle zinnen in~$\Gamma$. Om dit laatste te bewijzen,
hebben we nodig dat $\lnot !A \in \Gamma^*$ dan en slechts dan als
$!A \notin \Gamma^*$, dat $(!A \lor !B) \in \Gamma^*$ dan en slechts
dan als $!A \in \Gamma^*$ of $!B \in \Gamma^*$, enzovoort.
\end{explain}

In het vervolg zullen we de reflexiviteit, monotonie en transitiviteit
van~$\Proves$ vaak stilzwijgend gebruiken (zie
\iftag{FOL}{%
  \tagrefs{prfSC/{fol:seq:ptn:sec},prfND/{fol:ntd:ptn:sec},prfAX/{fol:axd:ptn:sec},prfTab/{fol:tab:ptn:sec}}}{%
  \tagrefs{prfSC/{pl:seq:ptn:sec},prfND/{pl:ntd:ptn:sec},prfAX/{pl:axd:ptn:sec},prfTab/{pl:tab:ptn:sec}}}).

\begin{prop}
\ollabel{prop:ccs}
Veronderstel dat $\Gamma$ !!{complete} en consistent is. Dan geldt:
\begin{enumerate}
\item \ollabel{prop:ccs-prov-in} Als $\Gamma \Proves !A$, dan $!A \in
  \Gamma$.

\tagitem{prvAnd}{\ollabel{prop:ccs-and} $!A \land !B \in \Gamma$
  dan en slechts dan als zowel $!A \in \Gamma$ als $!B \in \Gamma$.}{}

\tagitem{prvOr}{\ollabel{prop:ccs-or} $!A \lor !B \in \Gamma$ dan en
  slechts dan als $!A \in \Gamma$ of $!B \in \Gamma$.}{}

\tagitem{prvIf}{\ollabel{prop:ccs-if} $!A \lif !B \in \Gamma$ dan en
  slechts dan als $!A \notin \Gamma$ of $!B \in \Gamma$.}{}
\end{enumerate}
\end{prop}

\begin{proof}
In het gehele volgende bewijs nemen we aan dat $\Gamma$ !!{complete} en
consistent is.
\begin{enumerate}
\item Als $\Gamma \Proves !A$, dan $!A \in \Gamma$.

Veronderstel dat $\Gamma \Proves !A$. Neem ter tegenspraak aan dat $!A
\notin \Gamma$. Omdat $\Gamma$ !!{complete} is, geldt dan
$\lnot !A \in \Gamma$. Volgens
\iftag{FOL}{%
  \tagrefs{prfSC/{fol:seq:prv:prop:explicit-inc},
    prfND/{fol:ntd:prv:prop:explicit-inc},
    prfAX/{fol:axd:prv:prop:explicit-inc},
    prfTab/{fol:tab:prv:prop:explicit-inc}}}{%
  \tagrefs{prfSC/{pl:seq:prv:prop:explicit-inc},
    prfND/{pl:ntd:prv:prop:explicit-inc},
    prfAX/{pl:axd:prv:prop:explicit-inc},
    prfTab/{pl:tab:prv:prop:explicit-inc}}}
is $\Gamma$ dan inconsistent. Dit is in tegenspraak met de aanname dat
$\Gamma$ consistent is. Het kan dus niet zo zijn dat $!A \notin \Gamma$,
zodat $!A \in \Gamma$.

\tagitem{defAnd}{}{%
\iftag{probAnd}{Opgave.}{%
$!A \land !B \in \Gamma$ dan en slechts dan als zowel $!A \in \Gamma$
als $!B \in \Gamma$:

Voor de voorwaartse richting veronderstellen we dat
$!A \land !B \in \Gamma$. Dan volgt uit
\iftag{FOL}{%
  \tagrefs{prfSC/{fol:seq:ppr:prop:provability-land},
    prfND/{fol:ntd:ppr:prop:provability-land},
    prfAX/{fol:axd:ppr:prop:provability-land},
    prfTab/{fol:tab:ppr:prop:provability-land}}}{%
  \tagrefs{prfSC/{pl:seq:ppr:prop:provability-land},
    prfND/{pl:ntd:ppr:prop:provability-land},
    prfAX/{pl:axd:ppr:prop:provability-land},
    prfTab/{pl:tab:ppr:prop:provability-land}}}, onderdeel~(1),
dat $\Gamma \Proves !A$ en $\Gamma \Proves !B$. Volgens
\olref{prop:ccs-prov-in} geldt daarom $!A \in \Gamma$ en
$!B \in \Gamma$, zoals vereist.

Voor de omgekeerde richting nemen we $!A \in \Gamma$ en
$!B \in \Gamma$ aan. Uit
\iftag{FOL}{%
  \tagrefs{prfSC/{fol:seq:ppr:prop:provability-land},
    prfND/{fol:ntd:ppr:prop:provability-land},
    prfAX/{fol:axd:ppr:prop:provability-land},
    prfTab/{fol:tab:ppr:prop:provability-land}}}{%
  \tagrefs{prfSC/{pl:seq:ppr:prop:provability-land},
    prfND/{pl:ntd:ppr:prop:provability-land},
    prfAX/{pl:axd:ppr:prop:provability-land},
    prfTab/{pl:tab:ppr:prop:provability-land}}}, onderdeel~(2),
volgt dat $\Gamma \Proves !A \land !B$. Volgens
\olref{prop:ccs-prov-in} is dus $!A \land !B \in \Gamma$.}}

\tagitem{defOr}{}{%
\iftag{probOr}{Opgave.}{%
Eerst tonen we aan dat uit $!A \lor !B \in \Gamma$ volgt dat
$!A \in \Gamma$ of $!B \in \Gamma$. Veronderstel dat
$!A \lor !B \in \Gamma$, maar dat $!A \notin \Gamma$ en
$!B \notin \Gamma$. Omdat $\Gamma$ !!{complete} is, geldt
$\lnot !A \in \Gamma$ en $\lnot !B \in \Gamma$. Uit
\iftag{FOL}{%
  \tagrefs{prfSC/{fol:seq:ppr:prop:provability-lor},
    prfND/{fol:ntd:ppr:prop:provability-lor},
    prfAX/{fol:axd:ppr:prop:provability-lor},
    prfTab/{fol:tab:ppr:prop:provability-lor}}}{%
  \tagrefs{prfSC/{pl:seq:ppr:prop:provability-lor},
    prfND/{pl:ntd:ppr:prop:provability-lor},
    prfAX/{pl:axd:ppr:prop:provability-lor},
    prfTab/{pl:tab:ppr:prop:provability-lor}}},
onderdeel~(1), volgt dan dat $\Gamma$ inconsistent is, een tegenspraak.
Bijgevolg geldt $!A \in \Gamma$ of $!B \in \Gamma$.

Voor de omgekeerde richting veronderstellen we dat $!A \in \Gamma$ of
$!B \in \Gamma$. Uit
\iftag{FOL}{%
  \tagrefs{prfSC/{fol:seq:ppr:prop:provability-lor},
    prfND/{fol:ntd:ppr:prop:provability-lor},
    prfAX/{fol:axd:ppr:prop:provability-lor},
    prfTab/{fol:tab:ppr:prop:provability-lor}}}{%
  \tagrefs{prfSC/{pl:seq:ppr:prop:provability-lor},
    prfND/{pl:ntd:ppr:prop:provability-lor},
    prfAX/{pl:axd:ppr:prop:provability-lor},
    prfTab/{pl:tab:ppr:prop:provability-lor}}}, onderdeel~(2),
volgt dat $\Gamma \Proves !A \lor !B$. Volgens
\olref{prop:ccs-prov-in} is dus $!A \lor !B \in \Gamma$, zoals
vereist.}}

\tagitem{defIf}{}{%
\iftag{probIf}{Opgave.}{%
Voor de voorwaartse richting veronderstellen we dat
$!A \lif !B \in \Gamma$. Neem ter tegenspraak aan dat
$!A \in \Gamma$ en $!B \notin \Gamma$. Op grond van deze aannamen
gelden $!A \lif !B \in \Gamma$ en $!A \in \Gamma$. Uit
\iftag{FOL}{%
  \tagrefs{prfSC/{fol:seq:ppr:prop:provability-lif},
    prfND/{fol:ntd:ppr:prop:provability-lif},
    prfAX/{fol:axd:ppr:prop:provability-lif},
    prfTab/{fol:tab:ppr:prop:provability-lif}}}{%
  \tagrefs{prfSC/{pl:seq:ppr:prop:provability-lif},
    prfND/{pl:ntd:ppr:prop:provability-lif},
    prfAX/{pl:axd:ppr:prop:provability-lif},
    prfTab/{pl:tab:ppr:prop:provability-lif}}}, onderdeel~(1),
volgt dat $\Gamma \Proves !B$. Maar volgens
\olref{prop:ccs-prov-in} geldt dan $!B \in \Gamma$, in tegenspraak met
de aanname dat $!B \notin \Gamma$.

Voor de omgekeerde richting beschouwen we eerst het geval
$!A \notin \Gamma$. Omdat $\Gamma$ !!{complete} is, geldt
$\lnot !A \in \Gamma$. Uit
\iftag{FOL}{%
  \tagrefs{prfSC/{fol:seq:ppr:prop:provability-lif},
    prfND/{fol:ntd:ppr:prop:provability-lif},
    prfAX/{fol:axd:ppr:prop:provability-lif},
    prfTab/{fol:tab:ppr:prop:provability-lif}}}{%
  \tagrefs{prfSC/{pl:seq:ppr:prop:provability-lif},
    prfND/{pl:ntd:ppr:prop:provability-lif},
    prfAX/{pl:axd:ppr:prop:provability-lif},
    prfTab/{pl:tab:ppr:prop:provability-lif}}}, onderdeel~(2),
volgt dat $\Gamma \Proves !A \lif !B$. Opnieuw volgens
\olref{prop:ccs-prov-in} verkrijgen we $!A \lif !B \in \Gamma$, zoals
vereist.

Beschouw nu het geval $!B \in \Gamma$. Uit
\iftag{FOL}{%
  \tagrefs{prfSC/{fol:seq:ppr:prop:provability-lif},
    prfND/{fol:ntd:ppr:prop:provability-lif},
    prfTab/{fol:tab:ppr:prop:provability-lif},
    prfAX/{fol:axd:ppr:prop:provability-lif}}}{%
  \tagrefs{prfSC/{pl:seq:ppr:prop:provability-lif},
    prfND/{pl:ntd:ppr:prop:provability-lif},
    prfTab/{pl:tab:ppr:prop:provability-lif},
    prfAX/{pl:axd:ppr:prop:provability-lif}}},
opnieuw onderdeel~(2), volgt dat $\Gamma \Proves !A \lif !B$. Volgens
\olref{prop:ccs-prov-in} is dus $!A \lif !B \in \Gamma$.}}
\end{enumerate}
\end{proof}

\tagprob{FOL}
\begin{prob}
Voltooi het bewijs van \olref[fol][com][ccs]{prop:ccs}.
\end{prob}
\tagendprob

\tagprob{notFOL}
\begin{prob}
Voltooi het bewijs van \olref[pl][com][ccs]{prop:ccs}.
\end{prob}
\tagendprob

\end{document}
